\documentclass[UTF8]{ctexart}
\usepackage[paperwidth=17.6cm,paperheight=9.8cm,margin=0.08cm]{geometry}
\usepackage{tikz}
\usetikzlibrary{arrows.meta,positioning,fit,backgrounds,calc}
\pagestyle{empty}

\definecolor{mainblue}{RGB}{58,103,186}
\definecolor{deepblue}{RGB}{43,80,155}
\definecolor{goldline}{RGB}{181,143,35}
\definecolor{redline}{RGB}{214,72,67}
\definecolor{greenline}{RGB}{52,154,92}
\definecolor{purpleline}{RGB}{120,103,190}
\definecolor{grayline}{RGB}{110,110,110}

\tikzset{
  box/.style={
    rounded corners=5pt,
    draw=deepblue,
    line width=0.7pt,
    fill=mainblue,
    text=white,
    align=center,
    font=\small,
    inner sep=5pt,
    minimum height=1.18cm
  },
  smallbox/.style={
    box,
    font=\footnotesize,
    minimum height=0.88cm
  },
  groupbox/.style={
    draw=#1,
    dashed,
    line width=0.85pt,
    rounded corners=4pt,
    inner sep=7pt
  },
  flow/.style={
    -{Stealth[length=2.7mm,width=2.1mm]},
    draw=deepblue,
    line width=1.15pt
  },
  softflow/.style={
    -{Stealth[length=2.4mm,width=1.9mm]},
    draw=grayline,
    line width=0.8pt,
    dashed
  },
  title/.style={
    font=\small\bfseries,
    align=center,
    text=black
  },
  side/.style={
    font=\small\bfseries,
    align=center,
    text=black
  }
}

\begin{document}
\noindent
\begin{tikzpicture}[x=1cm,y=1cm,shift={(1.78,4.50)}]

  \node[title] at (0,4.15) {问题来源};
  \node[title] at (3.0,4.15) {理论基础};
  \node[title] at (6.15,4.15) {执行层控制器};
  \node[title] at (9.45,4.15) {参考层仲裁器};
  \node[title] at (12.75,4.15) {综合验证};

  \node[box, text width=2.20cm, minimum height=5.55cm] (ch1) at (0,0.55) {
    第一章\\
    绪论与问题凝练\\[0.25em]
    \footnotesize
    轨迹精度\\
    接触力安全\\
    人类在线修正\\
    几何边界约束
  };

  \node[box, text width=2.30cm, minimum height=5.55cm] (ch2) at (3.0,0.55) {
    第二章\\
    控制理论基础\\[0.25em]
    \footnotesize
    任务空间建模\\
    柔性接触模型\\
    阻抗与力位控制\\
    合作博弈理论
  };

  \node[box, text width=2.45cm] (ch3a) at (6.15,2.65) {
    第三章\\
    合作博弈\\
    力位协同控制
  };
  \node[smallbox, text width=2.45cm, below=0.23cm of ch3a] (ch3b) {
    力位增广状态\\
    $\xi$ 与接触力误差
  };
  \node[smallbox, text width=2.45cm, below=0.18cm of ch3b] (ch3c) {
    折扣帕累托代价\\
    反馈律与增益调度
  };
  \node[smallbox, text width=2.45cm, below=0.18cm of ch3c] (ch3d) {
    力安全裕度\\
    $\rho_F,\alpha_{\mathrm{FP}}$
  };

  \node[box, text width=2.45cm] (ch4a) at (9.45,2.65) {
    第四章\\
    人机动态仲裁
  };
  \node[smallbox, text width=2.45cm, below=0.23cm of ch4a] (ch4b) {
    操作者输入特征\\
    幅值、持续时间与方向
  };
  \node[smallbox, text width=2.45cm, below=0.18cm of ch4b] (ch4c) {
    安全参考投影\\
    切向修正与法向抑制
  };
  \node[smallbox, text width=2.45cm, below=0.18cm of ch4c] (ch4d) {
    参考层权重\\
    $\alpha_{\mathrm{HR}}$
  };

  \node[box, text width=2.50cm] (ch5a) at (12.75,1.80) {
    第五章\\
    综合实验验证
  };
  \node[smallbox, text width=2.50cm, below=0.23cm of ch5a] (ch5b) {
    无 RCM 柔性表面接触
  };
  \node[smallbox, text width=2.50cm, below=0.18cm of ch5b] (ch5c) {
    带 RCM 长工具操作
  };
  \node[smallbox, text width=2.50cm, below=0.18cm of ch5c] (ch5d) {
    控制器对比\\
    仲裁消融与实物验证
  };

  \node[box, fill=deepblue, text width=2.75cm, minimum height=0.95cm] (ch6) at (12.75,-3.18) {
    第六章\quad 总结与展望
  };

  \begin{pgfonlayer}{background}
    \node[groupbox=goldline, fit=(ch1)] (g1) {};
    \node[groupbox=grayline, fit=(ch2)] (g2) {};
    \node[groupbox=redline, fit=(ch3a)(ch3b)(ch3c)(ch3d)] (g3) {};
    \node[groupbox=greenline, fit=(ch4a)(ch4b)(ch4c)(ch4d)] (g4) {};
    \node[groupbox=purpleline, fit=(ch5a)(ch5b)(ch5c)(ch5d)] (g5) {};
  \end{pgfonlayer}

  \draw[flow] (ch1.east) -- (ch2.west);
  \draw[flow] (ch2.east) -- (ch3a.west);
  \draw[flow] (ch3a.east) -- (ch4a.west);
  \draw[flow] (ch4a.east) -- (ch5a.west);
  \draw[flow] (ch5d.south) -- (ch6.north);

  \draw[softflow] ($(ch3d.east)+(0.05,-0.02)$) .. controls (8.1,-1.70) and (10.55,-1.70) .. (ch5d.west);
  \draw[softflow] ($(ch4d.east)+(0.05,0.02)$) .. controls (10.90,-1.20) and (11.40,-1.15) .. (ch5c.west);

  \node[side, rotate=90] at (-1.65,0.55) {研究背景与科学问题};
  \node[side, rotate=90] at (14.55,-0.45) {平台验证与论文收束};

  \node[font=\footnotesize, align=center, text=black!75] at (6.15,-4.05) {
    主线关系：工程需求提出问题，理论基础支撑建模，第三章解决执行层力位协调，第四章解决参考层人机仲裁，第五章在统一平台验证二者组合效果
  };

\end{tikzpicture}
\end{document}
